% !TEX program = xelatex
% ============================================================
% DragonAI Beamer 模板骨架(starter)—— 复制本文件,替换内容即可开写
%
% 编译:xelatex starter.tex  ×2(第二遍生成目录、总页数)
% 交可编辑 PPT 时:把 \usetheme 行换成 \usetheme[ppt]{dragonai},
%   另存一份编译,拿转出的 PDF 去转 PPT(详见《使用说明.md》第 7 节)
%
% 原则:本文件里只有内容和结构命令,没有一句排版代码;
%   版式、间距、颜色、字体全部由主题接管,溢出了先删字再想别的。
% ============================================================
\documentclass[aspectratio=169,11pt,t]{beamer}
\usetheme{dragonai}

% ---------- 元信息:只改文字 ----------
\title{在这里写标题\\需要断行用双反斜杠}
\subtitle{在这里写副标题(不需要就整行删掉)}
\author{在这里写姓名}
\institute{单位(不需要就整行删掉)}
\date{2026.07}

\begin{document}

\DAcover     % 封面(全自动)
\DAtoc       % 目录(全自动;章节多时行距自动收紧)

% ============ 第 1 章 ============
\DAsection[一句给这章定调的话(红底语句,不需要可连方括号删掉)]{第一章标题}

% ---- 要点页:frame{标题}{副标题},副标题可省 ----
\begin{frame}{页标题}{页副标题(可删)}
  \begin{itemize}
    \item \textbf{要点一} —— 展开说明,行内强调用 \DAhot{红底强调}
    \item \textbf{要点二} —— 需要二级要点就再嵌一层
    \begin{itemize}
      \item 二级要点
    \end{itemize}
  \end{itemize}
  \DAtagline{落点}{一页的结论写在这里,比块版式省一半竖向空间}
  \DAtagline*{备注}{注解性的尾行用星号版,文字自动弱化为灰}
\end{frame}

% ---- 双栏 + 图:左栏说话,右栏放证据 ----
\begin{frame}{双栏页标题}
  \begin{columns}[T]
    \begin{column}{0.46\textwidth}
      \begin{itemize}
        \item 左栏要点一
        \item 左栏要点二
      \end{itemize}
      \medskip
      \DAstatbox{-92\%}{红底数字 + 一句说明}
    \end{column}
    \begin{column}{0.50\textwidth}
      \DAfigureph[37mm]{图占位:有真图后换成下一行的写法}
      % \DAfigbox{37mm}{你的图.png}{图注(注明出处)}
    \end{column}
  \end{columns}
\end{frame}

% ---- 图行页:多图并排,\hfill 均布;单图就只留一个,自动居中 ----
\begin{frame}{图行页标题}
  \begin{DAfigrow}
    \DAfigureph[26mm][40mm]{图一占位}\hfill
    \DAfigureph[26mm][40mm]{图二占位}
    % 真图写法(尺寸只给高度,宽度随图):
    % \DAfigbox{26mm}{a.png}{图注}\hfill \DAfigbox{26mm}{b.png}{图注}
  \end{DAfigrow}
  \DAtagline{看图说话}{图行下面配一句解读,别让听众自己猜}
\end{frame}

% ============ 第 2 章 ============
\DAsection{第二章标题}

% ---- 表格页:≤6 行用 DAtable;整页总表换 DAtable*(写法相同) ----
\begin{frame}{表格页标题}{整页总表时,把表注写在这个副标题里}
  \begin{DAtable}[表注(不需要可连方括号删掉)]{@{}lll@{}}
    \toprule
    列一 & 列二 & 列三 \\
    \midrule
    \DAtableno{01} & 第一行内容 & 说明 \\
    \DAtableno{02} & 第二行内容 & 说明 \\
    \bottomrule
  \end{DAtable}
\end{frame}

\DAclosing[谢谢]   % 结尾(联系行想换就在导言区 \renewcommand{\DAcontact}{…})

\end{document}
